% from ega4-18-fr.tex
\oldpage[IV-4]{197}

même de $f'$ (2.1.4), la réciproque étant vraie si $g$ est fidèlement plat ((2.5.1) et (2.7.1)); en outre pour tout $s'\in S'$, si $s=g(s')$, il est équivalent de dire que $f^{-1}(s)$ est une intersection complète ou que $f'^{-1}(s')$ est une intersection complète (19.3.4); d'où (ii), compte tenu de ce que $g$ est surjectif lorsqu'il est fidèlement plat.

Enfin, sous les hypothèses de (iii), $g\circ f$ est plat et localement de présentation finie. Par définition (19.3.6), on est donc ramené au cas où $Z$ est le spectre d'un corps, et à prouver alors que les anneaux locaux $\sh{O}_{X,x}$ sont des anneaux d'intersection complète, ce qui est contenu dans le résultat plus général suivant:

\begin{corollary}[19.3.10]
\label{IV.19.3.10}
Soient $f:X\to S$ un morphisme plat localement de présentation finie qui soit un morphisme d'intersection complète, $x$ un point de $X$, $s=f(x)$. Si $\sh{O}_{S,s}$ est un anneau noethérien d'intersection complète, il en est de même de $\sh{O}_{X,x}$.
\end{corollary}

On peut évidemment se borner au cas où $S=\Spec(A)$ avec $A=\sh{O}_{S,s}$. Montrons en outre qu'on peut se borner au cas où l'anneau local $A$ est \emph{complet}. En effet, posons $A'=\widehat A$, $X'=X\otimes_A A'$, et soit $s'$ l'unique point fermé de $S'=\Spec(A')$, qui est l'unique point de $S'$ au-dessus de $s$; si $f'=f_{(S')}:X'\to S'$, la fibre $f'^{-1}(s')$ est isomorphe canoniquement à $f^{-1}(s)$, puisque $A$ et $A'$ ont même corps résiduel (I, 3.6.4); il y a donc un seul point $x'\in X'$ au-dessus de $x$ et de $s'$, et on a par suite $\sh{O}_{X',x'}=\sh{O}_{X,x}\otimes_A A'$. On en déduit que les anneaux locaux $\sh{O}_{X,x}$ et $\sh{O}_{X',x'}$ ont des \emph{complétés isomorphes}, car de façon générale, si $E$ est un anneau local qui est une $A$-algèbre (l'homomorphisme $A\to E$ étant local), le séparé complété $(E\otimes_A\widehat A)^{\wedge}$ de l'anneau $E\otimes_A\widehat A$ muni de la topologie produit tensoriel est isomorphe au séparé complété de $\widehat E\widehat\otimes_A\widehat A\cong\widehat E$, donc au séparé complété $\widehat E$ de $E$ ($0_{\mathrm I}$, 7.7.1). En vertu de la définition (19.3.1), il revient donc au même de dire que $\sh{O}_{X,x}$ est un anneau d'intersection complète, ou que $\sh{O}_{X',x'}$ est un anneau d'intersection complète.

Supposons donc $A$ \emph{complet}; on peut en outre supposer que $X=\Spec(B)$, $B$ étant quotient d'une algèbre de polynômes $C=A[T_1,\ldots,T_r]$. Comme les anneaux de polynômes sur un corps sont des anneaux réguliers ($0$, 17.3.7), il résulte de (19.3.7) que l'immersion $X\to Z=\Spec(C)$ peut être supposée transversalement régulière relativement à $S$; donc on peut supposer que $B=C/\mathfrak J$, où $\mathfrak J$ est engendré par une suite régulière $(g_i)_{1\leq i\leq n}$ d'éléments de $C$. D'autre part, puisque $A$ est complet, on peut par hypothèse l'écrire $A'/\mathfrak R$, où $A'$ est un anneau local régulier et $\mathfrak R$ un idéal de $A'$ ($0$, 19.8.8), et l'hypothèse que $A$ est un anneau d'intersection complète entraîne que $\mathfrak R$ est engendré par une suite régulière $(f_j)_{1\leq j\leq m}$ (19.3.2). Dans l'anneau de polynômes $C'=A'[T_1,\ldots,T_r]$, les éléments $f_j$ ($1\leq j\leq m$) forment encore une suite régulière engendrant l'idéal $\mathfrak R'=\mathfrak R[T_1,\ldots,T_r]$ ($0$, 15.1.4); comme $C'/\mathfrak R'=C$, on voit que si, pour chaque $i$, $g'_i$ est un élément de $C'$ d'image $g_i$ dans $C$, la suite formée des $f_j$ ($1\leq j\leq m$) et des $g'_i$ ($1\leq i\leq n$) est régulière dans $C'$; si $\mathfrak J'$ est l'idéal de $C'$ qu'elle engendre, on a $B=C'/\mathfrak J'$, d'où la conclusion en vertu de (19.3.2), puisque $C'$ est un anneau régulier ($0$, 17.3.7).

\subsection{Application: critères de régularité et de lissité pour les préschémas éclatés.}
\EGAReaderTarget{ega4.c04.subsec.19-4}{subsection:IV.19.4}

